\chapter{The Original Method: Table-Specialist}\label{ch:paper}

This chapter explains the method of Xing et al.~\cite{xing2024tablellmspecialist} and its released code
\cite{tablespecialistrepo}. It covers how the method is meant to work, why each part is needed, and exactly which
details the paper and code leave unspecified, because those gaps shaped every later decision. The paper was first read
verbatim from its arXiv HTML version~2 for Sections~1--5, while its Section~6 tables for tasks other than NL-to-SQL could
at first be read only through a summary. On 2026-11-18 the published EMNLP 2025 version~\cite{xing2025emnlp} was read
in full, and every figure below was checked against its Table~3. Section and figure numbers below refer to arXiv
version~2 unless marked as the published version.
The mathematical notation below is ours; it formalises the paper's algorithms as the paper and repository describe them.

\section{Background and motivation}

A \emph{generalist} model such as an instruction-tuned LLM can be prompted for any table task. A \emph{specialist}
fine-tuned for one task is usually more accurate and cheaper to run, but needs labelled training data. Table tasks are
expensive to label, because a label is often a program (SQL, Scala, R), a cell, or a mapping between columns.
Table-Specialist replaces the human labeller with two models that check each other.

The key observation is that many table tasks come in \emph{dual} pairs:
\begin{itemize}
\item a \emph{generative} form, which produces an answer for an input, and
\item a \emph{classificative} or checking form, which judges whether an answer is right.
\end{itemize}
Sometimes the second form is simply a different program for the same question. If both forms can be run, their
agreement is evidence of correctness that needs no human label. The method then fine-tunes the generator on its own
validated outputs and repeats.

\section{Problem formulation}\label{sec:formulation}

Let $\mathcal{T}$ be a corpus of real, unlabelled tables, and let a task map an input $x$ (a table, a question) to an
output $y$ (a program, a set of cells). We want parameters $\theta$ of a model $M_\theta$ that maximise accuracy on
held-out benchmarks of that task, using only $\mathcal{T}$ and no labels.

\paragraph{Generator and validator.} A round of the method uses three components.
\begin{itemize}
\item A \emph{question model} $Q$ writes an input for a table: $x = Q(t)$, $t\in\mathcal{T}$.
\item A \emph{generator} $M_G$ writes the task output, $y = M_G(x)$.
\item A \emph{validator} $M_C$ writes the dual output $y' = M_C(x)$, or checks $y$ directly.
\end{itemize}
A validation function $V(x, y, y') \in \{0,1\}$ accepts or rejects the example. The round's training set is
\begin{equation}
  \mathcal{D}_r \;=\; \bigl\{\, (x, y) \;:\; t \in \mathcal{T}_r,\; x = Q(t),\; y = M_G^{(r-1)}(x),\;
  y' = M_C^{(r-1)}(x),\; V(x,y,y') = 1 \,\bigr\},
  \label{eq:dataset}
\end{equation}
where $\mathcal{T}_r$ is the round's slice of tables and $M^{(0)}$ is the vanilla model. Training on $\mathcal{D}_r$
gives $M^{(r)}$, which generates in round $r+1$. The paper runs ``up to 3'' rounds.

\paragraph{Execution validation (NL-to-SQL).} For NL-to-SQL, $y$ is SQL and $y'$ is a Spark-Scala program for the same
question. Let $S_1,\dots,S_K$ be random row samples of table $t$, each holding a fraction $\phi$ of the rows. Let
$\operatorname{run}(p, S)$ be the result of program $p$ on sample $S$. Then
\begin{equation}
  V_{\text{exec}}(x,y,y') \;=\; \prod_{k=1}^{K} \mathbb{1}\!\left[\operatorname{run}(y,S_k) \equiv
  \operatorname{run}(y',S_k)\right],
  \label{eq:vexec}
\end{equation}
with $K=10$ and $\phi=0.5$ in the repository's sampling script (\fpath{prepare_exec.py}). The paper's Algorithm~3
writes the test as ``$r \neq r'$''. The equivalence $\equiv$ is not further defined, which matters because SQL and
Spark name columns differently and return rows in different orders (\cref{sec:comparator}).

\paragraph{Permutation validation (error detection).} For error detection, the generator injects an error $e$ into a
column whose distinct values are $U$. The validator must find it regardless of the order in which the values are shown.
With permutations $\pi_1,\dots,\pi_J$ ($J=5$, seeds $0$--$4$ in the repository),
\begin{equation}
  V_{\text{perm}}(U,e) \;=\; \prod_{j=1}^{J} \mathbb{1}\!\left[\, M_C\bigl(\pi_j(U \cup \{e\})\bigr) = \{e\} \,\right],
  \qquad
  V_{\text{perm}}^{\text{clean}}(U) \;=\; \prod_{j=1}^{J} \mathbb{1}\!\left[\, M_C\bigl(\pi_j(U)\bigr) = \varnothing
  \,\right].
  \label{eq:vperm}
\end{equation}
A positive example passes if the validator names exactly the injected value under every order. A clean column passes
if it names nothing under every order. Kept positives are labelled with $e$, kept clean columns with ``NULL''.

\paragraph{Why each component is needed.}
\begin{itemize}
\item The \emph{dual check} is the only source of supervision; without it the model trains on its own mistakes.
The paper's no-validation ablation tests exactly this.
\item \emph{Sampling} several sub-tables (\cref{eq:vexec}) guards against chance agreement on one table. Two wrong
programs may agree on the full table, but rarely on ten random halves.
\item \emph{Permutation} (\cref{eq:vperm}) guards against position bias: a validator that always names the last value
would otherwise agree with an error placed last.
\item \emph{Iteration} lets a better model generate better data, and a better validator keep more of it.
\end{itemize}

\section{The algorithms}\label{sec:paper-algorithms}

\Cref{alg:paper-loop,alg:paper-exec,alg:paper-perm} restate the paper's Algorithms~1--3 in our notation. Comments mark
where a detail comes from the repository rather than the paper.

\begin{algorithm}[htbp]
\caption{Table-Specialist iterative generator--validator fine-tuning (paper Algorithm~1, restated)}\label{alg:paper-loop}
\KwIn{table corpus $\mathcal{T}$; vanilla models $M_G^{(0)}, M_C^{(0)}$; rounds $R$ (``up to 3''); tables per round $n$
(``e.g.\ 3000'')}
\KwOut{specialist $M_G^{(R)}$}
\For{$r \leftarrow 1$ \KwTo $R$}{
  $\mathcal{T}_r \leftarrow$ next $n$ tables of $\mathcal{T}$\Comment*[r]{repo: 1st, 2nd, 3rd 3K files}
  $\mathcal{D}^G_r \leftarrow \varnothing$;\ $\mathcal{D}^C_r \leftarrow \varnothing$\Comment*[r]{training sets reset each round}
  \ForEach{$t \in \mathcal{T}_r$}{
    $x \leftarrow Q(t)$\Comment*[r]{question model: not specified}
    $y \leftarrow M_G^{(r-1)}(x)$;\ $y' \leftarrow M_C^{(r-1)}(x)$\;
    \If{$V(x,y,y') = 1$}{
      add $(x,y)$ to $\mathcal{D}^G_r$;\ add $(x,y')$ to $\mathcal{D}^C_r$\;
    }
  }
  $M_G^{(r)} \leftarrow \operatorname{FineTune}(\,\cdot\,, \mathcal{D}^G_r)$;\ $M_C^{(r)} \leftarrow
  \operatorname{FineTune}(\,\cdot\,, \mathcal{D}^C_r)$\Comment*[r]{from previous model or from vanilla: ambiguous}
}
\Return{$M_G^{(R)}$}
\end{algorithm}

\begin{algorithm}[htbp]
\caption{Execution-based validation for NL-to-SQL (paper Algorithm~3, restated)}\label{alg:paper-exec}
\KwIn{table $t$; SQL $y$; dual program $y'$ (Spark-Scala); samples $K=10$; fraction $\phi=0.5$\Comment*[r]{repo}}
\KwOut{keep or reject}
\For{$k \leftarrow 0$ \KwTo $K-1$}{
  $S_k \leftarrow$ \texttt{t.sample(frac=}$\phi$\texttt{, random\_state=}$k$\texttt{)}\Comment*[r]{repo \fpath{prepare_exec.py}}
  $r \leftarrow \operatorname{run}_{\text{SQL}}(y, S_k)$;\ $r' \leftarrow \operatorname{run}_{\text{Spark}}(y', S_k)$\;
  \lIf{either run fails \textbf{or} $r \neq r'$}{\Return{reject}}
}
\Return{keep}
\end{algorithm}

\begin{algorithm}[htbp]
\caption{Permutation-based validation for error detection (paper Algorithm~2, restated)}\label{alg:paper-perm}
\KwIn{column values $U$; injected error $e$ (or none for a clean column); validator $M_C$; permutations $J=5$}
\KwOut{keep or reject}
$W \leftarrow U \cup \{e\}$ (or $U$ for a clean column)\;
\For{$j \leftarrow 0$ \KwTo $J-1$}{
  $W_j \leftarrow$ \texttt{random.seed(}$j$\texttt{); random.shuffle(}$W$\texttt{)}\Comment*[r]{repo \fpath{ed_validation_prompt_gen_v3.py}}
  $A \leftarrow$ values $M_C$ names as errors in $W_j$\;
  \lIf{$A \neq \{e\}$ (or $A \neq \varnothing$ for a clean column)}{\Return{reject}}
}
\Return{keep}
\end{algorithm}

\begin{figure}[htbp]
\centering
\begin{adjustbox}{max width=\linewidth}\begin{tikzpicture}[
  node distance=0.55cm and 0.7cm,
  box/.style={draw=navy,fill=white,rounded corners=3pt,thick,align=center,font=\sffamily\small,minimum height=1.0cm,
              text width=2.6cm,blur shadow={shadow blur steps=4,shadow xshift=0.6pt,shadow yshift=-0.6pt}},
  model/.style={box,fill=headshade},
  check/.style={box,draw=amber,fill=paperbg},
  data/.style={box,draw=sea,fill=oursbg},
  unspec/.style={box,dashed,draw=coral},
  arr/.style={-{Stealth[length=5pt]},thick,color=slate}]
  \node[data] (corpus) {Table corpus\\\footnotesize slice $\mathcal{T}_r$, 3,000 tables};
  \node[unspec,right=of corpus] (q) {Question model $Q$\\\footnotesize (unspecified)};
  \node[model,above right=0.25cm and 0.7cm of q] (g) {Generator $M_G^{(r-1)}$\\\footnotesize writes SQL};
  \node[model,below right=0.25cm and 0.7cm of q] (c) {Dual model $M_C^{(r-1)}$\\\footnotesize writes Spark-Scala};
  \node[check,right=2.9cm of q] (v) {Execute on 10 row\\samples; all agree?};
  \node[data,below=1.2cm of v] (d) {Training sets\\$\mathcal{D}^G_r$, $\mathcal{D}^C_r$};
  \node[unspec,left=of d] (ft) {LoRA fine-tune\\\footnotesize from vanilla or previous?};
  \node[model,left=of ft] (next) {$M_G^{(r)}$, $M_C^{(r)}$\\\footnotesize next round};
  \node[right=0.35cm of v,font=\sffamily\small,text=coral] (rej) {reject};
  \draw[arr] (corpus) -- (q);
  \draw[arr] (q) |- (g);
  \draw[arr] (q) |- (c);
  \draw[arr] (g) -| (v);
  \draw[arr] (c) -| (v);
  \draw[arr] (v) -- node[right,font=\sffamily\footnotesize]{agree} (d);
  \draw[arr,coral] (v) -- (rej);
  \draw[arr] (d) -- (ft);
  \draw[arr] (ft) -- (next);
  \draw[arr,dashed] (next.north) to[out=110,in=250] node[left,font=\sffamily\footnotesize,align=center]{round\\$r+1$} (corpus.south);
\end{tikzpicture}
\end{adjustbox}
\caption{The Table-Specialist loop for NL-to-SQL as described by the paper. One round turns a slice of unlabelled tables
into validated (question, SQL) and (question, Scala) pairs, and fine-tunes both models. The new models generate the next
round. Dashed boxes mark components the paper does not specify (who writes the questions; whether fine-tuning
restarts from vanilla).}
\label{fig:paper-loop}
\end{figure}

\section{Tasks, data and prompts}

\begin{paperbox}[title={Tasks with results in the paper (Llama-3.1-8B)}]
NL-to-SQL, NL-to-Scala, NL-to-R, schema matching, error detection, table question answering, and data
transformation (Pandas, R and SQL). Entity matching, column-type annotation and fact verification are listed in the
paper's Table~1 but have no results.
\end{paperbox}

\Cref{tab:paper-tasks} lists the test sets and validation per task (published version, Table~2 and Section~3).

{\small
\rowcolors{2}{rowshade}{white}
\begin{longtable}{@{}p{2.6cm}p{4.2cm}p{4.6cm}p{3.2cm}@{}}
\caption{Tasks in the paper, their dual-task validation, and their test sets. Sizes and metrics are from the published
Table~2.}\label{tab:paper-tasks}\\
\toprule\rowcolor{headshade}
Task & Validation (dual task) & Test sets (cases) & Metric\\\midrule\endfirsthead
\toprule\rowcolor{headshade}
Task & Validation (dual task) & Test sets (cases) & Metric\\\midrule\endhead
\bottomrule\endfoot
NL-to-SQL & SQL vs Spark-Scala execution agreement (\cref{eq:vexec}) & WikiSQL (1,000), WikiTQ (1,000),
Text2Analysis (271), Spider (1,198), BIRD (356) & Execution accuracy\\
NL-to-Scala & The same pairs; Scala checked against SQL (Example~5) & The same five & Execution accuracy\\
NL-to-R & Execution agreement with a second language (not named) & The same five & Execution accuracy\\
Schema matching & Model-based, consistent under permutation (Algorithm~2) & DeepM (42), WikiData (24), HXD (468) &
F1\\
Error detection & Permutation consistency of an injected error (\cref{eq:vperm}) & Spreadsheet-Tables (1,126),
Relational-Tables (1,081) & F1\\
Table QA & Paired with NL-to-code tasks; details not given & FinQA (1,000), TableBench (424), WikiTQ (1,000) & Accuracy\\
Data transformation & Execution agreement & TDE (570), Transform-Text (335) & Execution accuracy\\
\end{longtable}
}

\paragraph{Training tables.} The paper describes its corpus only as ``a large corpus'' of real tables. Its repository's
data file names (\fpath{table2question_all_webtable_v4_1st_3K}) suggest web tables, sliced into 3,000-table files per
round. The corpus is not released.

\paragraph{Prompts.} The repository releases the prompts as code.
\begin{itemize}
\item \textbf{Generation.} Few-shot scripts with four worked examples, such as
\fpath{Generation/NL2DSL/NL2SQL/nl2sql_generation_v6.py}. Tables are serialised as pipe-delimited markdown.
\item \textbf{Questions.} \fpath{table2question_generationv_v4.py} samples ``requirements'' for each question
(predicates 70\%, other requirements 33\% each).
\item \textbf{Errors.} Separate scripts write typo errors and semantic errors.
\item \textbf{Test and training format.} Prompt classes under \fpath{Finetune/data_generator/} build the zero-shot
prompts by default.
\end{itemize}
The only query path, \fpath{Evaluete/queryWithDeployment.py}, is a chat call with a fixed system message: ``You are an
AI assistant. Provide an answer to some data/table related questions.'' This is written for GPT-3.5; the paper does not
publish its Llama call.

\section{Training}\label{sec:paper-training}

\begin{paperbox}[title={Fine-tuning, as published}]
LoRA~\cite{hu2021lora}; batch size 1\% of the training data; a ``learning-rate multiplier of 0.5'' (an OpenAI
fine-tuning term). Section~6.5 and Figure~14 (published version: Appendix~B.2, Figure~10) compare restarting each round from the vanilla model with continuing from
the previous round, and find restarting ``consistently better''. The paper does not state which variant its main
results use, nor the LoRA rank, alpha, epochs or absolute learning rate.
\end{paperbox}

For a weight matrix $W_0\in\mathbb{R}^{d\times k}$, LoRA trains a low-rank update
\begin{equation}
  W \;=\; W_0 + \frac{\alpha}{r}\, B A, \qquad B\in\mathbb{R}^{d\times r},\; A\in\mathbb{R}^{r\times k},\; r \ll
  \min(d,k),
  \label{eq:lora}
\end{equation}
with $W_0$ frozen. Supervised fine-tuning minimises the negative log-likelihood of the target output given the prompt,
\begin{equation}
  \mathcal{L}(\theta) \;=\; -\frac{1}{|\mathcal{D}_r|}\sum_{(x,y)\in\mathcal{D}_r} \sum_{i=1}^{|y|}
  \log p_\theta\!\left(y_i \mid x, y_{<i}\right).
  \label{eq:sft}
\end{equation}
With a batch of 1\% of the data, every epoch has about 100 optimiser steps, whatever the size of $\mathcal{D}_r$.

\section{Evaluation}\label{sec:paper-metrics}

\paragraph{Execution accuracy.} For $N$ test cases with gold program $p_i^\ast$ and predicted program $\hat p_i$, run
on the case's table or database $D_i$,
\begin{equation}
  \mathrm{EA} \;=\; \frac{1}{N}\sum_{i=1}^{N} \mathbb{1}\!\left[\operatorname{norm}\bigl(\operatorname{run}(\hat
  p_i,D_i)\bigr) = \operatorname{norm}\bigl(\operatorname{run}(p_i^\ast,D_i)\bigr)\right].
  \label{eq:ea}
\end{equation}
In the repository's evaluators, $\operatorname{norm}$ converts values to numbers where possible and fills missing values
with the string ``None''. Equality is the list comparison \texttt{are\_lists\_equal}, which is order-sensitive and has a
relative tolerance for numbers. Its defects are analysed in \cref{sec:scoring}. For WikiTQ the gold answer is a list of
values, not a program.

\paragraph{Error-detection F1.} The repository's \fpath{ed_evaluator.py} computes precision only over columns with a
prediction, $\mathcal{P}=\{c:\hat E_c\neq\varnothing\}$, and recall only over columns with errors,
$\mathcal{R}=\{c:E_c\neq\varnothing\}$:
\begin{equation}
  P=\frac{1}{|\mathcal{P}|}\sum_{c\in\mathcal{P}}\frac{|\hat E_c\cap E_c|}{|\hat E_c|},\qquad
  R=\frac{1}{|\mathcal{R}|}\sum_{c\in\mathcal{R}}\frac{|\hat E_c\cap E_c|}{|E_c|},\qquad
  F_1=\frac{2PR}{P+R}.
  \label{eq:edf1}
\end{equation}
AutoTest~\cite{autotestrepo} also marks some values as ``debatable''. The repository ignores this field, and the paper
does not say how it was treated.

\section{Reported results}\label{sec:paper-results}

\Cref{tab:paper-results} gives the targets of this reproduction, from the published Table~3. The paper does not state
which test prompt or which checkpoint (base or instruction-tuned) its Llama column uses. For GPT-3.5, every Table~3 value
equals the last iteration plotted in Figure~6 (V3 for NL-to-code, V2 for error detection and schema matching), so Table~3
most likely reports the final round; no such figure is given for Llama.

{\small
\rowcolors{2}{rowshade}{white}
\begin{longtable}{@{}lrrrrrr@{}}
\caption{The paper's Llama-3.1-8B results that this report compares against (published Table~3). The mean of 4 excludes
Text2Analysis, which is unreleased; the paper's own averages include it.}\label{tab:paper-results}\\
\toprule\rowcolor{headshade}
Task and model & WikiSQL & WikiTQ & T2A & Spider & BIRD & Mean of 4\\\midrule\endfirsthead
\toprule\rowcolor{headshade}
Task and model & WikiSQL & WikiTQ & T2A & Spider & BIRD & Mean of 4\\\midrule\endhead
\bottomrule\endfoot
NL-to-SQL, vanilla & 0.525 & 0.300 & 0.273 & 0.670 & 0.388 & 0.4708\\
NL-to-SQL, Table-Specialist & 0.816 & 0.449 & 0.465 & 0.690 & 0.438 & 0.5983\\
NL-to-Scala, vanilla & 0.359 & 0.043 & 0.129 & 0.249 & 0.100 & 0.1878\\
NL-to-Scala, Table-Specialist & 0.728 & 0.188 & 0.214 & 0.356 & 0.228 & 0.3750\\
NL-to-R, vanilla & 0.331 & 0.138 & 0.103 & 0.503 & 0.225 & 0.2993\\
NL-to-R, Table-Specialist & 0.409 & 0.257 & 0.199 & 0.536 & 0.261 & 0.3658\\
\midrule
\multicolumn{7}{@{}l}{\textit{Other tasks}}\\
Error detection ST / RT, vanilla & \multicolumn{6}{l}{0.071 / 0.108 (F1)}\\
Error detection ST / RT, TS & \multicolumn{6}{l}{0.136 / 0.161 (F1)}\\
Schema matching, vanilla $\to$ TS & \multicolumn{6}{l}{0.839 $\to$ 0.873 (F1, average of 3 sets)}\\
Table QA, vanilla $\to$ TS & \multicolumn{6}{l}{0.270 $\to$ 0.296 (average of 3 sets, Table~4)}\\
\end{longtable}
}

Three points about the paper's own numbers, settled by the published version:
\begin{itemize}
\item \textbf{The ablation reports round~1.} The ablation table (arXiv Table~11, published Table~6, GPT-3.5 only) gives
NL-to-SQL 0.594, while Table~3 gives 0.616. This is not a contradiction: the ablation's Table-Specialist values (0.594,
0.496, 0.466 for SQL, R, Scala) equal the V1 values in Figure~6, and Table~3's equal V3. In that ablation a
language-model validator gives 0.595 against execution validation's 0.594, with vanilla at 0.569.
\item \textbf{The paper's averages include Text2Analysis.} The NL-to-Scala Llama average of 0.343 is the mean over five
sets, including Text2Analysis (0.214); over the four public sets it is 0.375. Likewise NL-to-R: 0.333 over five, 0.366
over four.
\item \textbf{No-validation numbers exist only for GPT-3.5}, as data labels in Figure~6 (NL-to-SQL 0.591, 0.591, 0.589
over three iterations). There are none for Llama.
\end{itemize}

\section{Assumptions and limitations of the original approach}\label{sec:paper-limits}

\begin{itemize}
\item \textbf{Agreement is not correctness.} Two programs can agree and both be wrong, for example when both misread a
vague question in the same way. Correct pairs can also be rejected because the two languages format results
differently. The paper reports no precision or recall for its validator, and no pass rates.
\item \textbf{Strict agreement discards most data.} Requiring all 10 samples to agree keeps a minority of pairs. In our
runs it kept 11--25\% (\cref{sec:sql-validation}). The paper discards rejected pairs with no repair or retry.
\item \textbf{The dual task must be executable or checkable.} Execution validation needs a second language with a
sandbox. Permutation validation relies on the validator being better than chance, and early rounds use the vanilla
model as validator.
\item \textbf{Self-reinforcement.} The validator in later rounds is trained on data it accepted itself. This may lock in
its biases, a risk the paper does not analyse.
\item \textbf{Under-specification.} The corpus, the question model, the restart choice for the main results, the LoRA
settings, the test prompt and the checkpoint are all unstated; the round reported is only implied by Figure~6
(\cref{tab:gaps}).
\end{itemize}

{\small
\rowcolors{2}{rowshade}{white}
\begin{longtable}{@{}p{4.1cm}p{5.4cm}p{5.2cm}@{}}
\caption{What the paper and repository specify, and what they leave open. Our choice for each gap is in
\cref{tab:adaptations}.}\label{tab:gaps}\\
\toprule\rowcolor{headshade}
Component & Specified & Not specified\\\midrule\endfirsthead
\toprule\rowcolor{headshade}
Component & Specified & Not specified\\\midrule\endhead
\bottomrule\endfoot
Model & ``Llama-3.1-8B''; README links \texttt{meta-llama/Llama-3.1-8B} (base) & Base or Instruct; revision\\
Query path & Repo: chat completions with a fixed system message & Any path for a model without a chat template; token
limits; temperature\\
Corpus & ``A large corpus'' of real tables; 3,000 per round ``e.g.'' & The tables; filters; pass rates; kept-pair
counts\\
Question model & Repo prompt and requirement sampler & Which model writes questions in later rounds\\
Fine-tuning & LoRA; batch 1\%; LR multiplier 0.5 & Rank; alpha; epochs; absolute LR; target modules\\
Rounds & Up to 3; restarting is better (Section 6.5) & Which variant Table~3 uses; the round only implied (Figure~6)\\
Validation & Algorithm 3; repo samples 10 $\times$ 50\% & How results are compared (order, names, types, tolerance)\\
Test prompt & Repo builds zero-shot test prompts by default & Which prompt produced Table 3\\
Test cases & Counts per test set & Case ID lists; WikiSQL/WikiTQ sampling; BIRD filter\\
Error detection & Algorithm 2; repo prompts and negatives & Who writes errors; validator per round; treatment of
debatable labels\\
NL-to-R & Repo generation prompt and evaluator & Any validation script\\
\end{longtable}
}
